\documentclass[11pt]{article}

\usepackage[margin=0.85in]{geometry}
\usepackage{amsmath, amssymb}
\usepackage{xcolor}
\usepackage{titlesec}
\usepackage{enumitem}
\usepackage{tcolorbox}

\definecolor{main}{HTML}{1E3A8A}
\definecolor{lightblue}{HTML}{EFF6FF}

\titleformat{\section}{\large\bfseries\color{main}}{}{0em}{}
\titleformat{\subsection}{\normalsize\bfseries}{}{0em}{}

\setlist[itemize]{noitemsep, topsep=2pt}

\tcbset{
    mybox/.style={
        colback=lightblue,
        colframe=main,
        arc=6pt,
        boxrule=0.8pt,
        left=8pt,
        right=8pt,
        top=8pt,
        bottom=8pt
    }
}

\begin{document}

\begin{center}
{\LARGE \textbf{AP Statistics: Unit 1}} \\[4pt]
{\small \textit{Statistics is the disciplined extraction of signal from noise.}} \\
{\small \textbf{~ By Lostless1907}}
\end{center}

\vspace{10pt}


\begin{tcolorbox}[mybox]
\section*{Population, Sample, and Units}

\textbf{Population:} complete collection of interest  

\textbf{Sample:} subset used for study  

\textbf{Statistical Unit:} a single element  

\textbf{Subject:} a human unit  

\subsection*{From Whole to Part}

A dataset is represented as:
\[
x_1, x_2, \dots, x_n
\]

Each $x_i$ corresponds to one unit.  
In practice, we observe $S \subseteq P$ and attempt to infer properties of $P$.

The quality of conclusions depends critically on how $S$ is chosen.
\end{tcolorbox}


\begin{tcolorbox}[mybox]
\section*{Parameters and Statistics}

\textbf{Parameter:} numerical summary of a population  

\textbf{Statistic:} numerical summary of a sample  

\[
\mu = \text{population mean}, \quad 
\bar{x} = \frac{1}{n}\sum_{i=1}^{n} x_i
\]

\subsection*{Estimation as the Core Problem}

Parameters are fixed but unknown.  
Statistics are random variables depending on the sample.

All of statistics is built on approximating parameters using statistics, and understanding the error in that approximation.
\end{tcolorbox}


\begin{tcolorbox}[mybox]
\section*{Variables}

\textbf{Variable:} a measurable attribute  

\textbf{Explanatory Variable ($x$):} predictor  

\textbf{Response Variable ($y$):} outcome  

\subsection*{Structure Behind Data}

A fundamental model:
\[
y = f(x) + \varepsilon
\]

where $\varepsilon$ captures randomness, noise, and unobserved influences.

Even when a relationship exists, it is rarely perfectly deterministic.
\end{tcolorbox}


\begin{tcolorbox}[mybox]
\section*{Confounding}

A \textbf{confounding variable} affects both explanatory and response variables.

\subsection*{When Relationships Mislead}

If $Z$ influences both $X$ and $Y$, then:
\[
P(Y|X) \neq P(Y|X,Z)
\]

The observed association between $X$ and $Y$ may be partially or entirely due to $Z$, rather than a direct relationship.

This is the primary obstacle to causal inference.
\end{tcolorbox}


\begin{tcolorbox}[mybox]
\section*{Correlation and Causation}

\textbf{Correlation:} statistical association  

\textbf{Causation:} direct influence  

\subsection*{Sources of Association}

An observed relationship may arise from:
\begin{itemize}
\item $X \rightarrow Y$ (causal effect)
\item $Z \rightarrow X, Y$ (confounding)
\item random variation (coincidence)
\end{itemize}

Establishing causation requires controlled experimental structure, not just observed association.
\end{tcolorbox}


\begin{tcolorbox}[mybox]
\section*{Types of Studies}

\textbf{Observational Study:} no intervention  

\textbf{Experiment:} active assignment of treatments  

\subsection*{Seeing vs Controlling}

Observational studies measure existing conditions and can reveal patterns.

Experiments manipulate variables and, when properly designed, allow isolation of causal effects.
\end{tcolorbox}


\begin{tcolorbox}[mybox]
\section*{Experiments}

\textbf{Treatment:} imposed condition  

\textbf{Experimental Unit:} object receiving treatment  

\textbf{Response:} measured outcome  

\subsection*{Structure of an Experiment}

An experiment assigns treatments to units and measures responses, aiming to compare outcomes across different conditions.

The validity of conclusions depends on design, not just data.
\end{tcolorbox}


\begin{tcolorbox}[mybox]
\section*{Control}

\textbf{Control Group:} baseline for comparison  

\textbf{Control Variable:} factor held constant  

\subsection*{Isolating Effects}

Control removes alternative explanations.

Without a control group, differences in outcomes cannot be attributed confidently to the treatment.
\end{tcolorbox}


\begin{tcolorbox}[mybox]
\section*{Randomization}

Assignment of treatments using chance.

\subsection*{Eliminating Hidden Bias}

Randomization distributes both known and unknown confounders across groups.

This ensures that differences in outcomes are due to treatment rather than pre-existing differences.
\end{tcolorbox}


\begin{tcolorbox}[mybox]
\section*{Experimental Designs}

\textbf{Completely Randomized Design:} assign treatments purely at random  

\textbf{Randomized Block Design:} group similar units, randomize within groups  

\textbf{Matched Pairs Design:} compare highly similar units or repeated measures  

\subsection*{Reducing Unnecessary Variation}

Blocking and matching reduce variability from known sources, increasing the precision of comparisons.

Less noise leads to clearer detection of treatment effects.
\end{tcolorbox}


\begin{tcolorbox}[mybox]
\section*{Blinding}

\textbf{Single-Blind:} subjects unaware  

\textbf{Double-Blind:} subjects and evaluators unaware  

\textbf{Placebo:} inert treatment  

\subsection*{Preventing Behavioral Bias}

Awareness of treatment can influence both participant behavior and measurement.

Blinding removes expectation effects and preserves objectivity.
\end{tcolorbox}


\begin{tcolorbox}[mybox]
\section*{Replication}

Repeating an experiment across multiple units.

\subsection*{Stability Through Repetition}

Replication reduces the effect of random variation.

Reliable conclusions require consistent results across many observations.
\end{tcolorbox}


\begin{tcolorbox}[mybox]
\section*{Law of Large Numbers}

As sample size increases, sample statistics approach population parameters.

\subsection*{Convergence Through Size}

\[
\bar{x} \to \mu \quad \text{as } n \to \infty
\]

Larger samples reduce variability and improve accuracy of estimates.

This is the mathematical foundation of statistical reliability.
\end{tcolorbox}


\begin{tcolorbox}[mybox]
\section*{Census}

A \textbf{census} collects data from the entire population.

\subsection*{Limits of Completeness}

Although exact, censuses are often impractical due to cost, time, and logistical constraints.

Carefully designed samples usually provide sufficient accuracy with far greater efficiency.
\end{tcolorbox}

\end{document}